% AUTO-GENERATED by writeup/make_rehaul_ablations.py -- do not edit; re-run the script.
% tab:rh_ablation_macro (2026-10-29 addendum L; (iii) 2026-10-30): each macro release calendar added alone to
% HAR (Panel A, vs HAR) and to HAR + RM + FOMC (Panel B, vs HAR + RM + FOMC), then all six; every bar, the
% scored days without / with the row's release; DM on daily average QLIKE.  Calendar: src/data/loading.py macro_release_days.
\begin{table}[htbp]
\centering
\scriptsize
\setlength{\tabcolsep}{1.5pt}
\begin{threeparttable}
\caption{Macroeconomic release calendars by day type: change in QLIKE (\%)}
\label{tab:rh_ablation_macro}
\begin{tabular*}{\linewidth}{@{\extracolsep{\fill}}lrrrrrrr@{}}
\toprule
 & OLS & Ridge & Lasso & Elastic net & LightGBM & XGBoost & Forest \\
\midrule
\multicolumn{8}{@{}l}{\textit{A. HAR benchmark plus the calendar, against the HAR benchmark ($-$: the calendar adds accuracy)}} \\
\multicolumn{8}{@{}l}{Every bar} \\
\quad CPI & -- & 0.0 (0.1) & 0.0 (0.0) & 0.0 ($-$0.2) & 0.0 (0.0) & $-$0.3 ($-$1.9) & 0.2 (0.7) \\
\quad Employment & -- & $-$0.1 ($-$1.5) & $-$0.1 ($-$1.3) & $-$0.1 ($-$2.3) & $-$0.6 ($-$3.4) & $-$0.5 ($-$3.0) & $-$0.2 ($-$0.7) \\
\quad PPI & -- & 0.1 (2.0) & 0.1 (2.0) & 0.1 (1.7) & 0.0 (0.2) & $-$0.1 ($-$0.8) & 0.1 (0.6) \\
\quad GDP (advance) & -- & $-$0.1 ($-$1.0) & $-$0.1 ($-$1.0) & $-$0.1 ($-$0.9) & $-$0.5 ($-$3.0) & $-$0.6 ($-$3.5) & $-$0.3 ($-$0.9) \\
\quad Trade & -- & $-$0.1 ($-$1.2) & $-$0.1 ($-$1.2) & $-$0.1 ($-$1.5) & 0.0 ($-$0.1) & $-$0.2 ($-$1.8) & 0.4 (1.5) \\
\quad Retail sales & -- & 0.0 (0.4) & 0.1 (0.5) & 0.0 (0.3) & $-$0.2 ($-$1.0) & $-$0.4 ($-$2.7) & 0.1 (0.6) \\
\addlinespace[1pt]
\quad All \rhabNMacReleases{} (macro) & -- & 0.1 (0.3) & 0.1 (0.3) & 0.0 ($-$0.2) & $-$0.3 ($-$1.4) & $-$0.6 ($-$2.5) & 0.2 (0.6) \\
\addlinespace[2pt]
\multicolumn{8}{@{}l}{Non-release days} \\
\quad CPI & -- & 0.0 ($-$0.2) & 0.0 ($-$0.2) & 0.0 ($-$0.8) & 0.0 (0.1) & $-$0.3 ($-$1.9) & 0.2 (0.5) \\
\quad Employment & -- & $-$0.1 ($-$1.2) & $-$0.1 ($-$1.0) & $-$0.1 ($-$1.9) & $-$0.6 ($-$3.4) & $-$0.5 ($-$2.9) & $-$0.2 ($-$0.5) \\
\quad PPI & -- & 0.1 (1.5) & 0.1 (1.5) & 0.1 (1.3) & 0.1 (0.3) & $-$0.1 ($-$0.4) & 0.2 (0.9) \\
\quad GDP (advance) & -- & $-$0.1 ($-$0.6) & $-$0.1 ($-$0.6) & 0.0 ($-$0.5) & $-$0.5 ($-$2.8) & $-$0.6 ($-$3.3) & $-$0.2 ($-$0.7) \\
\quad Trade & -- & $-$0.2 ($-$1.6) & $-$0.1 ($-$1.6) & $-$0.1 ($-$1.9) & $-$0.1 ($-$0.4) & $-$0.3 ($-$2.2) & 0.3 (1.1) \\
\quad Retail sales & -- & 0.0 (0.4) & 0.0 (0.4) & 0.0 (0.1) & $-$0.2 ($-$0.9) & $-$0.4 ($-$2.7) & 0.1 (0.7) \\
\addlinespace[1pt]
\quad All \rhabNMacReleases{} (macro) & -- & 0.1 (0.3) & 0.1 (0.3) & 0.0 ($-$0.3) & $-$0.3 ($-$0.9) & $-$0.6 ($-$2.2) & 0.4 (0.9) \\
\addlinespace[2pt]
\multicolumn{8}{@{}l}{Release days} \\
\quad CPI & -- & 0.3 (0.9) & 0.3 (0.6) & 0.5 (2.4) & $-$0.3 ($-$0.5) & $-$0.3 ($-$0.5) & 0.8 (1.0) \\
\quad Employment & -- & $-$0.8 ($-$0.8) & $-$0.7 ($-$0.7) & $-$0.9 ($-$1.1) & $-$0.3 ($-$0.4) & $-$0.7 ($-$0.8) & $-$1.3 ($-$1.3) \\
\quad PPI & -- & 0.7 (1.6) & 0.7 (1.6) & 0.5 (1.5) & $-$0.6 ($-$1.3) & $-$1.0 ($-$1.9) & $-$1.1 ($-$1.4) \\
\quad GDP (advance) & -- & $-$2.6 ($-$1.7) & $-$3.1 ($-$1.6) & $-$1.7 ($-$1.7) & $-$2.6 ($-$1.7) & $-$3.1 ($-$2.2) & $-$4.2 ($-$2.2) \\
\quad Trade & -- & 0.7 (1.2) & 0.7 (1.2) & 0.5 (1.0) & 0.8 (1.3) & 1.3 (1.9) & 2.4 (2.7) \\
\quad Retail sales & -- & 0.1 (0.3) & 0.2 (0.4) & 0.4 (1.0) & $-$0.3 ($-$0.5) & 0.0 (0.1) & $-$0.1 ($-$0.1) \\
\addlinespace[1pt]
\quad All \rhabNMacReleases{} (macro) & -- & 0.1 (0.2) & 0.1 (0.2) & 0.0 (0.2) & $-$0.6 ($-$1.4) & $-$0.6 ($-$1.4) & $-$0.4 ($-$0.8) \\
\bottomrule
\end{tabular*}
\begin{tablenotes}[flushleft]
\footnotesize
\item \textit{Sign:} each cell is $100\,(L_{\text{row}}/L_{\text{ref}} - 1)$, the percentage change in mean QLIKE of the row's forecast against its reference, with the Diebold--Mariano statistic in parentheses (same sign; daily average QLIKE).  Panel~A, the row's release calendar added to the HAR benchmark set (reference: that set), as the FOMC row of Panel~B of Table~\ref{tab:rh_ablation}: negative, the loss falls when the calendar is added, i.e.\ the calendar adds accuracy; positive, it hurts.
\item \textit{Notes:} The release calendars are built in the same way as the FOMC calendar.  Rows: CPI and PPI, the consumer and producer price indexes; Employment, the employment situation; GDP (advance), the advance estimate of gross domestic product; Trade, international trade.  Macro: the \rhabNMacReleases{} calendars together (Panel~F of Table~\ref{tab:rh_ablation}).  Every bar: the scored bars of Table~\ref{tab:rh_oos} (2014--2024).  Release days are the scored days on which the row's release is published (macro: at least one of the releases), and non-release days are the other scored days.  Non-release days and release days split the scored days as the ordinary days and the FOMC statement days do for the FOMC calendar.  Release days: CPI \rhabNRelDaysCpi{}, Employment \rhabNRelDaysEmp{}, PPI \rhabNRelDaysPpi{}, GDP (advance) \rhabNRelDaysGdp{}, Trade \rhabNRelDaysTrade{}, Retail sales \rhabNRelDaysRetail{}; macro \rhabNRelDays{}.  FOMC statement days among them: CPI \rhabNRelFomcDaysCpi{}, Employment \rhabNRelFomcDaysEmp{}, PPI \rhabNRelFomcDaysPpi{}, GDP (advance) \rhabNRelFomcDaysGdp{}, Trade \rhabNRelFomcDaysTrade{}, Retail sales \rhabNRelFomcDaysRetail{}; macro \rhabNRelFomcDays{}.  The test runs on daily average QLIKE.  The models with a calendar keep the full set's frozen hyperparameters (FOMC-patched design); the HAR benchmark keeps those of Table~\ref{tab:rh_oos}.  OLS is not estimated on the HAR benchmark set plus a calendar.
\end{tablenotes}
\end{threeparttable}
\end{table}
\begin{table}[htbp]
\centering
\scriptsize
\setlength{\tabcolsep}{1.5pt}
\begin{threeparttable}
\ContinuedFloat
\caption{Macroeconomic release calendars by day type: change in QLIKE (\%) (continued)}
\begin{tabular*}{\linewidth}{@{\extracolsep{\fill}}lrrrrrrr@{}}
\toprule
 & OLS & Ridge & Lasso & Elastic net & LightGBM & XGBoost & Forest \\
\midrule
\multicolumn{8}{@{}l}{\textit{B. HAR + RM + FOMC plus the calendar, against HAR + RM + FOMC ($-$: the calendar adds accuracy)}} \\
\multicolumn{8}{@{}l}{Every bar} \\
\quad CPI & 0.0 (0.4) & 0.0 (0.0) & 0.0 (0.4) & 0.0 ($-$0.4) & 0.0 (0.0) & 0.0 (0.1) & 0.2 (2.4) \\
\quad Employment & 0.0 (0.1) & $-$0.1 ($-$1.0) & $-$0.1 ($-$0.5) & $-$0.1 ($-$1.0) & $-$0.1 ($-$1.4) & $-$0.1 ($-$0.8) & 0.0 (0.1) \\
\quad PPI & 0.0 (0.3) & 0.0 (0.3) & 0.0 (0.2) & 0.0 (0.1) & 0.0 (0.7) & 0.0 (0.2) & 0.2 (2.6) \\
\quad GDP (advance) & 0.1 (0.7) & 0.0 (0.3) & 0.1 (0.9) & 0.0 (0.0) & 0.0 ($-$0.6) & $-$0.1 ($-$0.9) & 0.0 (0.7) \\
\quad Trade & 0.0 (0.0) & $-$0.1 ($-$0.7) & 0.0 (0.0) & $-$0.1 ($-$1.1) & 0.0 (0.0) & $-$0.1 ($-$0.7) & 0.2 (3.2) \\
\quad Retail sales & $-$0.1 ($-$0.6) & $-$0.1 ($-$0.7) & $-$0.1 ($-$0.9) & $-$0.1 ($-$0.8) & 0.0 ($-$0.1) & 0.0 ($-$0.7) & 0.2 (2.7) \\
\addlinespace[1pt]
\quad All \rhabNMacReleases{} (macro) & 3.7 (1.4) & 0.0 (0.1) & 0.1 (0.4) & $-$0.1 ($-$0.5) & 0.0 (0.0) & $-$0.1 ($-$1.1) & 0.5 (3.4) \\
\addlinespace[2pt]
\multicolumn{8}{@{}l}{Non-release days} \\
\quad CPI & 0.0 ($-$0.5) & $-$0.1 ($-$0.8) & 0.0 ($-$0.5) & $-$0.1 ($-$1.3) & 0.0 ($-$0.5) & 0.0 ($-$0.1) & 0.2 (2.2) \\
\quad Employment & 0.0 (0.0) & $-$0.1 ($-$1.1) & $-$0.1 ($-$0.7) & $-$0.1 ($-$1.1) & $-$0.1 ($-$1.5) & 0.0 ($-$0.6) & 0.0 (0.2) \\
\quad PPI & 0.0 ($-$0.2) & 0.0 ($-$0.2) & 0.0 ($-$0.1) & 0.0 ($-$0.2) & 0.1 (0.8) & 0.0 (0.3) & 0.2 (2.8) \\
\quad GDP (advance) & 0.1 (0.9) & 0.1 (0.9) & 0.1 (1.4) & 0.0 (0.5) & 0.0 ($-$0.5) & $-$0.1 ($-$0.9) & 0.1 (0.9) \\
\quad Trade & $-$0.2 ($-$1.6) & $-$0.1 ($-$1.0) & 0.0 ($-$0.3) & $-$0.1 ($-$1.3) & 0.0 ($-$0.2) & 0.0 ($-$0.5) & 0.2 (2.7) \\
\quad Retail sales & $-$0.2 ($-$1.4) & $-$0.1 ($-$0.9) & $-$0.1 ($-$1.1) & $-$0.1 ($-$1.0) & 0.0 ($-$0.4) & $-$0.1 ($-$0.7) & 0.2 (2.6) \\
\addlinespace[1pt]
\quad All \rhabNMacReleases{} (macro) & 4.1 (1.2) & 0.0 (0.1) & 0.1 (0.2) & $-$0.1 ($-$0.6) & 0.0 ($-$0.4) & 0.0 ($-$0.3) & 0.5 (2.8) \\
\addlinespace[2pt]
\multicolumn{8}{@{}l}{Release days} \\
\quad CPI & 1.3 (3.0) & 0.9 (2.2) & 1.1 (2.8) & 0.9 (2.6) & 0.6 (2.2) & 0.3 (0.9) & 0.4 (0.9) \\
\quad Employment & 0.9 (1.1) & 0.5 (0.7) & 0.6 (0.8) & 0.1 (0.3) & 0.2 (0.6) & $-$0.5 ($-$1.5) & $-$0.4 ($-$1.3) \\
\quad PPI & 1.3 (1.8) & 0.6 (1.4) & 0.5 (1.4) & 0.4 (1.3) & 0.0 (0.0) & $-$0.1 ($-$0.2) & $-$0.1 ($-$0.5) \\
\quad GDP (advance) & $-$1.2 ($-$1.0) & $-$2.3 ($-$1.9) & $-$2.3 ($-$2.3) & $-$1.4 ($-$2.0) & $-$0.3 ($-$0.9) & $-$0.2 ($-$1.0) & $-$0.6 ($-$1.4) \\
\quad Trade & 3.7 (1.4) & 0.6 (2.0) & 0.6 (2.2) & 0.4 (1.4) & 0.2 (1.0) & $-$0.4 ($-$0.7) & 0.9 (2.7) \\
\quad Retail sales & 2.5 (1.3) & 0.2 (0.6) & 0.5 (1.1) & 0.4 (1.3) & 0.4 (1.2) & 0.0 ($-$0.1) & 0.2 (0.8) \\
\addlinespace[1pt]
\quad All \rhabNMacReleases{} (macro) & 2.3 (1.7) & 0.0 (0.0) & 0.2 (0.6) & 0.0 (0.0) & 0.2 (0.9) & $-$0.3 ($-$2.0) & 0.4 (2.1) \\
\bottomrule
\end{tabular*}
\begin{tablenotes}[flushleft]
\footnotesize
\item \textit{Sign:} as in Panel~A.  Panel~B, the row's release calendar added to HAR + RM + FOMC (reference: HAR + RM + FOMC): negative, the calendar adds accuracy; positive, it hurts.
\item \textit{Notes:} HAR + RM + FOMC: the HAR benchmark set plus the futures' realized moments (RM) and the FOMC calendar.  Calendars, days and test as in Panel~A.  The models with a calendar and HAR + RM + FOMC keep the full set's frozen hyperparameters (FOMC-patched design).
\end{tablenotes}
\end{threeparttable}
\end{table}
